\documentclass[12pt]{article}
\usepackage[T1]{fontenc}
\usepackage[utf8]{inputenc}
\usepackage[polish]{babel}
\usepackage{lmodern}
\usepackage[a4paper,margin=2.4cm]{geometry}
\usepackage{xcolor}
\usepackage{graphicx}
\usepackage{amsmath,amssymb}
\usepackage{booktabs}
\usepackage{array}
\usepackage{enumitem}
\usepackage{hyperref}
\usepackage{tcolorbox}
\usepackage{fancyhdr}
\definecolor{WEIIPrimary}{RGB}{0,130,150}
\definecolor{WEIIPrimaryDark}{RGB}{0,86,100}
\definecolor{WEIIPrimaryLight}{RGB}{229,246,248}
\definecolor{PRzNavy}{RGB}{31,63,142}
\definecolor{RLTextDark}{RGB}{34,34,34}
\definecolor{RLGray}{RGB}{90,90,90}
\definecolor{RLLightGray}{RGB}{245,246,248}
\definecolor{RLWarning}{RGB}{184,86,0}
\definecolor{RLExample}{RGB}{52,112,64}
\definecolor{RLCodeBg}{RGB}{248,248,248}

\usepackage{listings}
\usepackage{xcolor}

\lstdefinestyle{rlpython}{%
    language=Python,
    basicstyle=\ttfamily\scriptsize,
    keywordstyle=\color{WEIIPrimaryDark}\bfseries,
    commentstyle=\color{RLGray}\itshape,
    stringstyle=\color{RLExample},
    numbers=left,
    numberstyle=\tiny\color{RLGray},
    numbersep=6pt,
    showstringspaces=false,
    breaklines=true,
    tabsize=4,
    frame=single,
    rulecolor=\color{WEIIPrimary!35},
    backgroundcolor=\color{RLCodeBg},
    captionpos=b
}

\lstset{style=rlpython}


\hypersetup{colorlinks=true,linkcolor=WEIIPrimaryDark,urlcolor=WEIIPrimaryDark,citecolor=WEIIPrimaryDark}
\setlength{\headheight}{15pt}
\pagestyle{fancy}
\fancyhf{}
\lhead{Uczenie się ze wzmocnieniem}
\rhead{WEiI PRz}
\cfoot{\thepage}

\newtcolorbox{taskbox}[1]{colback=WEIIPrimaryLight,colframe=WEIIPrimary,title=\textbf{#1}}
\newtcolorbox{notebox}[1]{colback=RLLightGray,colframe=RLGray,title=\textbf{#1}}
\newtcolorbox{warninghandout}[1]{colback=RLWarning!8,colframe=RLWarning,title=\textbf{#1}}

\newcommand{\handouttitle}[3]{%
\begin{center}
{\Large\bfseries #1}\\[0.2cm]
{\large #2}\\[0.15cm]
{\small #3}
\end{center}
\vspace{0.4cm}
\hrule
\vspace{0.6cm}
}

\newcommand{\EnableLabNumbering}{%
  \renewcommand{\thesection}{L\arabic{section}}%
  \renewcommand{\thesubsection}{\thesection.\arabic{subsection}}%
  \renewcommand{\thesubsubsection}{\thesubsection.\arabic{subsubsection}}%
}


\title{Raport projektowy\\Uczenie się ze wzmocnieniem}
\author{Jakub Fuhrman}
\date{2026/2027}

\begin{document}
\maketitle
\tableofcontents
\newpage

\section{Informacje podstawowe}

\begin{center}
\begin{tabular}{p{0.34\textwidth}p{0.56\textwidth}}
\toprule
Element & Informacja \\
\midrule
Imię i nazwisko & Jakub Fuhrman \\
Numer albumu & 179482 \\
Kierunek/specjalność & Informatyka / Sztuczna inteligencja \\
Tytuł projektu & Agent-kurier na gridzie z energią i zamykającymi się przejściami \\
Link do repozytorium & \url{https://github.com/JakubFuh/UszW-kurier-grid} \\
Data oddania & 12.11.2026 \\
\bottomrule
\end{tabular}
\end{center}

\section{Temat i cel projektu}

\subsection{Opis problemu}

Projekt dotyczy problemu decyzyjnego w dyskretnym środowisku siatkowym (grid).
Agent wciela się w rolę kuriera poruszającego się po mapie z przeszkodami i pułapkami.
Jego zadaniem jest zbieranie paczek oraz dostarczanie ich do bazy.

Agent dysponuje ograniczonym zasobem energii, która maleje wraz z wykonywanymi ruchami.
Dodatkowo środowisko zawiera zdarzenia dynamiczne: po zebraniu określonej liczby paczek
wybrane przejścia na mapie zostają zamknięte. Wymusza to planowanie z wyprzedzeniem,
a nie wyłącznie lokalne dążenie do najbliższej nagrody.

W projekcie przewidziano kilka map o różnym poziomie trudności
(różny układ przeszkód, rozmieszczenie paczek, rozmieszczenie pułapek oraz moment zamykania przejść).

\subsection{Cel agenta}

Celem agenta jest dostarczenie wszystkich paczek do bazy przed wyczerpaniem energii
i przed przekroczeniem limitu kroków, unikając jednocześnie pułapek oraz zablokowania
na mapie po zamknięciu przejść.

Działanie agenta uznaje się za skuteczne, gdy w epizodzie:
\begin{itemize}
    \item wszystkie paczki zostaną zebrane i dostarczone do bazy,
    \item agent nie wejdzie w pułapkę,
    \item energia agenta nie spadnie do zera przed ukończeniem zadania,
    \item nie zostanie przekroczony limit kroków epizodu.
\end{itemize}

Dodatkowymi miarami jakości będą m.in. średni zwrot (return), skuteczność osiągania celu
oraz średnia długość epizodu.

\subsection{Uzasadnienie zastosowania uczenia się ze wzmocnieniem}

Problem nadaje się do uczenia się ze wzmocnieniem, ponieważ:
\begin{itemize}
    \item decyzje są podejmowane sekwencyjnie --- w każdym kroku agent wybiera kierunek ruchu;
    \item decyzje wpływają na przyszły stan środowiska --- zmieniają pozycję agenta, poziom energii,
    liczbę zebranych paczek oraz (po osiągnięciu progu) dostępność przejść na mapie;
    \item istotny jest wynik długoterminowy --- zebranie najbliższej paczki może później uniemożliwić
    powrót do bazy po zamknięciu przejść lub doprowadzić do wyczerpania energii;
    \item zastosowanie agenta uczącego się jest uzasadnione, ponieważ optymalna strategia zależy
    od układu mapy, kolejności zbierania paczek i zarządzania energią, a nie da się jej w prosty sposób
    sprowadzić do jednorazowej klasyfikacji lub statycznej optymalizacji.
\end{itemize}

\section{Formalizacja problemu RL/MDP}

\subsection{Agent i środowisko}

\textbf{Agent:} kurier poruszający się po dyskretnej siatce.\\
\textbf{Środowisko:} mapa gridowa zawierająca pola wolne, przeszkody, pułapki, paczki, bazę
oraz przejścia, które mogą zostać zamknięte w trakcie epizodu.

Interakcja w jednym kroku epizodu przebiega następująco:
\begin{enumerate}
    \item agent obserwuje bieżący stan (lub obserwację),
    \item wybiera akcję ruchu,
    \item środowisko aktualizuje pozycję agenta, energię oraz ewentualnie stan paczek/przejść,
    \item środowisko zwraca nagrodę $r_t$ oraz informację o zakończeniu epizodu.
\end{enumerate}

\subsection{Stan lub obserwacja}

W projekcie przyjmuje się pełną obserwowalność stanu (MDP).
Agent ma dostęp do informacji wystarczających do określenia aktualnej sytuacji na mapie.

\begin{center}
\begin{tabular}{p{0.30\textwidth}p{0.60\textwidth}}
\toprule
Składowa stanu/obserwacji & Znaczenie \\
\midrule
Pozycja agenta $(x, y)$ & Aktualne położenie kuriera na siatce \\
Energia $e$ & Pozostały zasób energii agenta \\
Maska / zbiór paczek & Informacja, które paczki zostały już zebrane \\
Flaga zamknięcia przejść & Czy zdarzenie zamknięcia przejść już wystąpiło \\
Pozycja bazy & Położenie punktu dostawy (stałe dla danej mapy) \\
\bottomrule
\end{tabular}
\end{center}

W implementacji stan będzie reprezentowany w formie dyskretnej, dogodnej dla algorytmu
tablicowego (Q-learning). Szczegóły kodowania stanu mogą zostać doprecyzowane po pierwszych testach środowiska.

\subsection{Akcje}

Zbiór akcji jest dyskretny i obejmuje ruchy ortogonalne po siatce.
Próba wejścia w przeszkodę, poza mapę albo w zamknięte przejście nie zmienia pozycji agenta
(akcja jest nieskuteczna), ale zużywa energię i generuje karę za krok.

\begin{center}
\begin{tabular}{p{0.30\textwidth}p{0.60\textwidth}}
\toprule
Akcja & Opis \\
\midrule
\texttt{UP} & Ruch o jedno pole w górę \\
\texttt{DOWN} & Ruch o jedno pole w dół \\
\texttt{LEFT} & Ruch o jedno pole w lewo \\
\texttt{RIGHT} & Ruch o jedno pole w prawo \\
\bottomrule
\end{tabular}
\end{center}

\subsection{Funkcja nagrody}

Wstępna funkcja nagrody (do weryfikacji eksperymentalnej):

\begin{equation}
    r_t =
    \begin{cases}
        +10 & \text{zebranie paczki}, \\
        +50 & \text{dostarczenie wszystkich paczek do bazy (sukces)}, \\
        -20 & \text{wejście w pułapkę}, \\
        -20 & \text{wyczerpanie energii}, \\
        -1 & \text{każdy wykonany krok}, \\
        0 & \text{w pozostałych przypadkach}.
    \end{cases}
\end{equation}

Uzasadnienie:
\begin{itemize}
    \item nagroda dodatnia za zebranie paczki motywuje do realizacji zadania cząstkowego,
    \item duża nagroda końcowa za sukces premiuje ukończenie całego epizodu,
    \item kary za pułapkę i wyczerpanie energii zniechęcają do ryzykownych strategii,
    \item kara za krok zachęca do krótszych tras i lepszego zarządzania energią,
    \item wartości nagród będą ewentualnie korygowane, jeśli okaże się, że agent uczy się
    niepożądanego zachowania (np. zbiera paczki, ale nie wraca do bazy).
\end{itemize}

\subsection{Warunki zakończenia epizodu}

Epizod kończy się, gdy zachodzi którakolwiek z sytuacji:
\begin{itemize}
    \item sukces --- wszystkie paczki zostały dostarczone do bazy,
    \item porażka --- agent wszedł w pułapkę,
    \item porażka --- energia agenta spadła do zera,
    \item porażka --- przekroczono maksymalną liczbę kroków w epizodzie.
\end{itemize}

\subsection{Założenia i ograniczenia środowiska}

Na obecnym etapie przyjmuje się następujące założenia:
\begin{itemize}
    \item środowisko jest w pełni obserwowalne,
    \item mapa jest dyskretna i skończona,
    \item przeszkody są statyczne; dynamiczne są wyłącznie zamykane przejścia,
    \item energia maleje o stałą wartość przy każdym kroku,
    \item nie modeluje się fizyki ciągłej ani częściowej obserwowalności (POMDP),
    \item szczegóły parametrów (rozmiar map, próg zamykania przejść, limity energii)
    zostaną doprecyzowane podczas implementacji środowiska.
\end{itemize}

Planowany algorytm uczenia: \textbf{Q-learning}, z możliwością późniejszego porównania z SARSA.
Jako agentów bazowych przewiduje się agenta losowego oraz agenta zachłannego
(np. idź w stronę najbliższej paczki / bazy).

\section{Implementacja}

\textit{Sekcja do uzupełnienia w kolejnych etapach projektu.}

\subsection{Struktura repozytorium}

\begin{verbatim}
UszW-kurier-grid/
|-- src/
|-- experiments/
|-- report/
|-- results/
|-- ai_transcripts/
|-- requirements.txt
`-- README.md
\end{verbatim}

\subsection{Środowisko}

Do uzupełnienia po implementacji środowiska (\texttt{reset}, \texttt{step}).

\subsection{Agent bazowy}

Do uzupełnienia: agent losowy oraz agent zachłanny.

\subsection{Agent uczący się}

Do uzupełnienia: Q-learning (ewentualnie porównanie z SARSA).

\section{Plan eksperymentów}

\textit{Sekcja do uzupełnienia w kolejnych etapach projektu.}

\begin{center}
\begin{tabular}{p{0.24\textwidth}p{0.30\textwidth}p{0.30\textwidth}}
\toprule
Eksperyment & Co jest zmieniane? & Co jest mierzone? \\
\midrule
E1 & Do uzupełnienia & Do uzupełnienia \\
E2 & Do uzupełnienia & Do uzupełnienia \\
E3 & Do uzupełnienia & Do uzupełnienia \\
\bottomrule
\end{tabular}
\end{center}

\section{Wyniki}

\textit{Sekcja do uzupełnienia po przeprowadzeniu eksperymentów.}

\section{Analiza i wnioski}

\textit{Sekcja do uzupełnienia po uzyskaniu wyników.}

\section{Instrukcja uruchomienia}

\textit{Sekcja do uzupełnienia po przygotowaniu uruchamialnej wersji projektu.}

\section{Wsparcie AI}

Na etapie przygotowania propozycji tematu i wstępnej formalizacji korzystano ze wsparcia AI
(Cursor / Composer) do konsultacji pomysłu projektowego oraz redakcji treści rozdziałów 1--3.
Transkrypty rozmów będą przechowywane w katalogu \texttt{ai\_transcripts/} repozytorium.
Szczegółowa tabela użycia AI będzie uzupełniana na bieżąco w trakcie realizacji projektu.

\begin{center}
\begin{tabular}{p{0.13\textwidth}p{0.19\textwidth}p{0.21\textwidth}p{0.27\textwidth}}
\toprule
Data & Narzędzie/model & Cel użycia & Transkrypt / link / zakres \\
\midrule
06--12.11.2026 & Cursor / Composer & Konsultacja tematu, formalizacja RL, redakcja raportu & czat w Cursor; katalog \texttt{ai\_transcripts/} \\
\bottomrule
\end{tabular}
\end{center}

Treść merytoryczna została sprawdzona i zaakceptowana przez autora projektu.
Ostateczna odpowiedzialność za poprawność formalizacji i późniejszą implementację
spoczywa na studencie.

\section{Literatura}

\begin{thebibliography}{9}
\bibitem{sutton_barto}
R. S. Sutton, A. G. Barto, \textit{Reinforcement Learning: An Introduction}, 2nd ed., MIT Press, 2018.
\bibitem{zasady_projektu}
Materiały przedmiotu \textit{Uczenie się ze wzmocnieniem}, WEiI PRz, zasady realizacji projektu, 2026/2027.
\end{thebibliography}

\end{document}
